\usepackage{ctex}
\usepackage{amsfonts,amsmath,amscd,amssymb,amsthm} %美国数学会套装
\usepackage{latexsym,amsbsy,bbold, fullpage}
\usepackage{pdfsync,yfonts}
\usepackage[all,pdf]{xy}
\usepackage{tkz-graph,tikz-cd,xy}
\usepackage{color,xcolor,tcolorbox,framed}      %颜色
\usepackage{fancyhdr,lastpage}                  %页眉页脚，最后一页
\usepackage{incgraph}
\usepackage{multicol}                           %双栏
\usepackage{enumerate}                          %枚举

\usepackage{geometry}
\newgeometry{left = 3.5 cm, right=3.5cm}

\usepackage[T1]{fontenc}                       %Required for accented characters
\usepackage{titlesec}                          %章节标题风格
\usepackage{hyperref,url}                      %超链接
\usepackage{mathpazo}                          %曲豆豆喜欢用的Palatino字体
\usepackage[framemethod=tikz]{mdframed}        %花里胡哨的框框
\usepackage{wallpaper}                         %背景图片
\usepackage{pstricks,pst-node,pst-tree}        %图论（暂时没卵用）
\usepackage{graphicx,wrapfig,float}            % 用来插图
\usepackage{ifthen}                            %逻辑判断
\usepackage{tikz,stmaryrd}                     %究极绘图语言
\usepackage{booktabs}                          %调整表格线与上下内容的间隔
\usepackage{chngpage}
\usepackage{enumitem}
\usepackage{enumerate}
\usepackage{makeidx}
\usepackage[tikz]{bclogo}

\usepackage{microtype}
%\usepackage[american]{babel}

\linespread{1.2} % 调整行间距
\allowdisplaybreaks[4]%允许跨页公式
%————————————————————————————————————————————————————

\newenvironment{solution}
  {\begin{proof}[解]\,}
  {\end{proof}}

%————————————————————————————————————————————————————

\newenvironment{prob}
  {\item\bfseries\kaishu}
  {}

\renewcommand{\proofname}{证明}

%-------------------------------------------------------------

\newcommand*{\Gray}[1]{\textcolor[rgb]{0.62,0.59,0.56}{#1}}

%——————————————————————————————————————————————————————————————————

\everymath{\displaystyle}  %全部行间公式，暗黑魔咒，慎用！

\renewcommand{\labelenumi}{\textbf{\theenumi}.\,}  %枚举环境编号设置
%\setlist[enumerate,2]{label=\arabic}
\renewcommand{\labelenumii}{({\theenumii})\,}  %枚举环境编号设置

%——————————索引标签术——————————————————%

\newcommand*{\newdef}[2]{
    \textbf{#1}(#2)
    \index{#2 $\quad$ #1}
    \label{def-#1}
  }

%——————————Hint———————————————————%

\newcommand*{\hint}[1]
  {
{\fangsong\textcolor[rgb]{0.6,0.6,0.6}{[提示:#1]}}
  }

\newcommand*{\remark}[1]
  {
    ({\fangsong #1})
  } 
